% ─────────────────────────────────────────────────────────────────────────────
% WERSJA POLSKA — ROZDZIAŁ 1 — WERSJE ROBOCZE
% Wszystkie wersje mają identyczną strukturę klasycznego wstępu:
%   1. Definicja i motywacja SLAM
%   2. Trudność problemu
%   3. Ewolucja historyczna
%   4. Luka badawcza
%   5. Cel i zakres pracy
%   6. Struktura pracy
% Różnią się wyłącznie stylem i rejestrem językowym.
% ─────────────────────────────────────────────────────────────────────────────


% ═════════════════════════════════════════════════════════════════════════════
% WERSJA A — formalny, klasyczny styl akademicki
% ═════════════════════════════════════════════════════════════════════════════

\chapter{Wstęp}

Jednym z~fundamentalnych problemów robotyki mobilnej, sformalizowanym w~latach
dziewięćdziesiątych XX~wieku, jest jednoczesna lokalizacja i~mapowanie
(ang.~\textit{Simultaneous Localisation and Mapping}, SLAM)~\cite{leonard1991mobile}.
Problem ten dotyczy zdolności robota do budowania spójnej mapy nieznanego środowiska
przy jednoczesnym wyznaczaniu własnej pozycji w~tej mapie.
SLAM stanowi warunek konieczny autonomicznej nawigacji bez sygnału GPS,
a~jego rozwiązanie jest niezbędne w~szerokim spektrum zastosowań ---
od pojazdów logistyki magazynowej i~dronów inspekcyjnych po roboty usługowe
działające w~przestrzeniach publicznych.

Trudność problemu wynika z~wzajemnej zależności jego dwóch składowych.
Dokładne budowanie mapy wymaga znajomości aktualnej pozycji robota,
natomiast precyzyjna lokalizacja wymaga spójnej mapy, względem której
pozycja ta może być wyznaczona.
Rozwiązywanie obu zadań jednocześnie, w~czasie rzeczywistym i~na podstawie
zaszumionych danych sensorycznych, pozostaje wyzwaniem algorytmicznym
i~obliczeniowym nawet po kilku dekadach intensywnych badań.

Od momentu formalnego sformułowania problemu dziedzina przeszła znaczącą ewolucję.
Pierwsze podejścia oparte na rozszerzonym filtrze Kalmana udowodniły
wykonalność koncepcji, lecz ich kwadratowa złożoność obliczeniowa ograniczała
zastosowanie do małych środowisk.
Zastosowanie filtrów cząsteczkowych, w~szczególności filtra Rao-Blackwellized,
umożliwiło obsługę większych map przy zachowaniu rozsądnego kosztu obliczeniowego.
Przełomem okazało się jednak grafowe sformułowanie SLAM, w~którym trajektoria
robota i~wiązania sensoryczne reprezentowane są jako graf podlegający globalnej
optymalizacji --- podejście to dominuje w~nowoczesnych implementacjach.
Równolegle rozszerzała się gama obsługiwanych sensorów: od dwuwymiarowych
skanerów laserowych, przez trójwymiarowy LiDAR, do kamer i~układów z~ciasno
sprzężoną centralą inercyjną (IMU).

Efektem tych postępów jest dziś bogaty ekosystem implementacji open-source,
lecz ich obiektywne porównanie pozostaje utrudnione.
Opublikowane benchmarki dotyczą zazwyczaj jednej rodziny metod lub jednego
środowiska i~rzadko przekładają się na odmienne konfiguracje sprzętowe.
Brakuje ujednoliconego zestawienia obejmującego różne modalności sensoryczne
i~rodziny algorytmów testowanych w~identycznych warunkach, co sprawia,
że dobór właściwego algorytmu dla konkretnej platformy sprzętowej pozostaje
zadaniem obarczonym znaczną niepewnością.

\section{Cel i~zakres pracy}

Celem niniejszej pracy jest systematyczna ocena i~porównanie algorytmów SLAM
różnych rodzin metod i~modalności sensorycznych, przeprowadzone w~jednolitych,
w~pełni kontrolowanych warunkach eksperymentalnych.
Ocenę przeprowadzono na robocie mobilnym TurtleBot3~Burger wyposażonym
w~trójwymiarowy skaner LiDAR, kamerę stereo, centralę inercyjną IMU
oraz enkodery odometrii kołowej.
Do badań wybrano jedenaście algorytmów obejmujących fuzję sensoryczną,
odometrię LiDAR, SLAM~2D i~3D na bazie LiDAR-u, ciasno sprzężoną odometrię
i~SLAM LiDAR-IMU oraz wizyjny SLAM.

Eksperymenty przeprowadzono na jednym zbiorze danych zarejestrowanym
w~symulowanym środowisku magazynowym w~NVIDIA~Isaac~Sim,
który dostarcza dokładnych pozycji referencyjnych (\textit{ground truth})
do obiektywnej ewaluacji.
Wspólnym środowiskiem uruchomieniowym dla wszystkich algorytmów był
ROS\,2 (\textit{Robot Operating System~2}).

Dokładność trajektorii oceniana jest miarami ATE i~RPE
odniesionymi do pozycji referencyjnych symulatora.
Jakość map 3D wyznaczana jest odległością Chamfera i~miarą F-score,
a~koszt obliczeniowy --- średnim obciążeniem procesora i~szczytowym
zużyciem pamięci RAM.

\section{Struktura pracy}

\begingroup
\setlength{\parskip}{2pt}
\setlength{\parindent}{1.5em}

Rozdział~2 zawiera podstawy teoretyczne niezbędne do zrozumienia omawianych
algorytmów, obejmując nawigację w~robotyce mobilnej, sensory, reprezentacje map,
rejestrację chmur punktów, filtracyjną i~grafową estymację stanu, preintegrację
IMU, podstawy wizyjnego SLAM oraz metryki ewaluacji.

Rozdział~3 opisuje każdy z~jedenastu ocenianych algorytmów pogrupowanych
według rodziny metod.

Rozdział~4 przedstawia konfigurację eksperymentalną, zbiór danych oraz wyniki
wszystkich jedenastu algorytmów w~zakresie dokładności trajektorii, jakości
mapy i~kosztu obliczeniowego.

Rozdział~5 podsumowuje wyniki eksperymentów, formułuje zalecenia dotyczące
doboru algorytmu SLAM do strukturyzowanych środowisk wewnętrznych oraz wskazuje
kierunki dalszych badań.

\endgroup


% ═════════════════════════════════════════════════════════════════════════════
% WERSJA B — bezpośredni, rzeczowy, technicznie precyzyjny
% ═════════════════════════════════════════════════════════════════════════════

\chapter{Wstęp}

Autonomiczny robot mobilny poruszający się bez sygnału GPS musi rozwiązywać dwa
zadania jednocześnie: budować mapę nieznanego środowiska i~wyznaczać w~niej
własną pozycję.
Problem ten, znany jako jednoczesna lokalizacja i~mapowanie
(ang.~\textit{Simultaneous Localisation and Mapping},
SLAM)~\cite{leonard1991mobile}, jest kluczowym elementem systemów nawigacji
robotów w~magazynach, szpitalach i~innych środowiskach pozbawionych
zewnętrznej infrastruktury pozycjonowania.

Zasadnicza trudność SLAM polega na tym, że oba zadania są od siebie zależne:
mapa jest potrzebna do lokalizacji, a~lokalizacja --- do budowy mapy.
Rozwiązanie tego problemu w~czasie rzeczywistym, przy zaszumionych danych
sensorycznych i~ograniczonych zasobach obliczeniowych robota,
od dekad stanowi aktywny obszar badań.

Podejścia algorytmiczne ewoluowały wyraźnie w~ciągu ostatnich trzydziestu lat.
Rozszerzony filtr Kalmana był pierwszą praktyczną metodą, lecz rósł kwadratowo
z~liczbą elementów mapy, co ograniczało jego zastosowanie.
Filtry cząsteczkowe poprawiły skalowalność i~umożliwiły obsługę nieliniowych
modeli ruchu.
Dominującym podejściem stało się grafowe SLAM, w~którym pozycje robota
i~wiązania sensoryczne tworzą graf optymalizowany globalnie --- umożliwia ono
efektywne wykrywanie i~zamykanie pętli nawet na długich trajektoriach.
Obsługiwane sensory rozszerzyły się od płaskich skanerów 2D do LiDAR-u 3D,
kamer stereo i~układów LiDAR-IMU ze sprzężeniem ciasnym.

Pomimo szerokiej dostępności implementacji open-source ich rzetelne porównanie
jest trudne --- istniejące benchmarki obejmują zwykle jedną rodzinę algorytmów
lub jeden typ sensora i~nie przekładają się na inne konfiguracje sprzętowe.
Brakuje jednego zestawienia, które pozwoliłoby inżynierowi wybrać algorytm SLAM
na podstawie pomiarów uzyskanych w~tych samych warunkach.

\section{Cel i~zakres pracy}

Celem pracy jest ocena i~porównanie jedenastu algorytmów SLAM różnych rodzin
i~modalności sensorycznych w~identycznych, kontrolowanych warunkach.
Eksperymenty przeprowadzono na robocie TurtleBot3~Burger (LiDAR 3D, kamera stereo,
IMU, enkodery kół).
Oceniane algorytmy obejmują fuzję sensoryczną, odometrię LiDAR, SLAM~2D i~3D
na bazie LiDAR-u, ciasno sprzężoną odometrię i~SLAM LiDAR-IMU oraz wizyjny SLAM.

Dane pochodzą z~symulowanego magazynu w~NVIDIA~Isaac~Sim, który dostarcza
dokładnych pozycji referencyjnych i~gwarantuje powtarzalność warunków.
Wszystkie algorytmy uruchomiono w~środowisku ROS\,2.

Trajektorie oceniane są miarami ATE i~RPE, mapy 3D --- odległością Chamfera
i~F-score, a~koszt obliczeniowy --- średnim obciążeniem CPU i~szczytowym
zużyciem RAM.

\section{Struktura pracy}

\begingroup
\setlength{\parskip}{2pt}
\setlength{\parindent}{1.5em}

Rozdział~2 zawiera podstawy teoretyczne: nawigację w~robotyce mobilnej, sensory,
reprezentacje map, rejestrację chmur punktów, filtracyjną i~grafową estymację
stanu, preintegrację IMU, podstawy wizyjnego SLAM oraz metryki ewaluacji.

Rozdział~3 opisuje każdy z~jedenastu ocenianych algorytmów pogrupowanych
według rodziny metod.

Rozdział~4 przedstawia konfigurację eksperymentalną, zbiór danych oraz wyniki
wszystkich jedenastu algorytmów: dokładność trajektorii, jakość mapy
i~koszt obliczeniowy.

Rozdział~5 podsumowuje wyniki, formułuje zalecenia dotyczące doboru algorytmu SLAM
do strukturyzowanych środowisk wewnętrznych i~wskazuje kierunki dalszych badań.

\endgroup


% ═════════════════════════════════════════════════════════════════════════════
% WERSJA C — płynna narracja, zdania złożone, styl eseistyczny
% ═════════════════════════════════════════════════════════════════════════════

\chapter{Wstęp}

Aby robot mobilny mógł poruszać się autonomicznie w~nieznanym środowisku,
musi równocześnie orientować się w~przestrzeni i~budować jej model ---
bez żadnej zewnętrznej infrastruktury, bez sygnału GPS, bez wcześniej
przygotowanej mapy.
To pozornie proste wymaganie kryje w~sobie jeden z~najtrudniejszych problemów
robotyki mobilnej, sformalizowany w~latach dziewięćdziesiątych XX~wieku
jako jednoczesna lokalizacja i~mapowanie
(ang.~\textit{Simultaneous Localisation and Mapping},
SLAM)~\cite{leonard1991mobile}.
Znalazło ono zastosowanie wszędzie tam, gdzie robot działa samodzielnie:
w~magazynach, szpitalach, na placach budowy i~w~obszarach niedostępnych dla ludzi.

U~podstaw trudności SLAM leży błędne koło: precyzyjna mapa środowiska wymaga
dokładnych pozycji robota w~trakcie jej tworzenia, a~dokładna lokalizacja robota
wymaga spójnej mapy, względem której pozycja jest wyznaczana.
Rozwiązanie tego wzajemnego uwikłania w~czasie rzeczywistym, przy nieuchronnym
szumie sensorycznym i~ograniczonej mocy obliczeniowej pokładowego komputera,
nie jest trywialne i~mimo kilku dekad intensywnych badań nadal pozostaje
aktywnym obszarem pracy naukowej.

Historia algorytmów SLAM jest historią kolejnych przełamywań ograniczeń.
Rozszerzony filtr Kalmana, dominujący w~pierwszej fazie badań, dostarczał
spójnych estymacji, lecz jego złożoność rosła kwadratowo z~liczbą elementów mapy,
co czyniło go niepraktycznym w~dużych środowiskach.
Filtry cząsteczkowe przyniosły lepszą skalowalność i~naturalną obsługę
wielomodalnych rozkładów prawdopodobieństwa, stając się podstawą szeroko stosowanych
systemów dwuwymiarowego SLAM.
Grafowe sformułowanie problemu --- w~którym trajektoria i~wiązania sensoryczne
tworzą sieć węzłów i~krawędzi podlegającą globalnej optymalizacji ---
umożliwiło efektywne działanie na długich trajektoriach i~skuteczne zamykanie pętli,
stając się dominującym paradygmatem nowoczesnych implementacji.
Towarzyszył temu równoległy rozwój bazy sensorycznej: od płaskich skanerów
laserowych 2D, przez bogate geometrycznie chmury punktów trójwymiarowego LiDAR-u,
do kamer i~układów z~ciasno sprzężoną centralą inercyjną, które poprawiają
odporność na dynamikę ruchu i~pozwalają korygować zniekształcenia pomiarów.

Bogactwo dostępnych implementacji open-source rodzi jednak praktyczny problem.
Każdy algorytm zaprojektowany jest pod konkretną kombinację sensorów, środowisk
i~wymagań obliczeniowych, a~opublikowane benchmarki obejmują zwykle wąski wycinek
tej przestrzeni --- jedną rodzinę metod, jeden typ sensora lub jedno środowisko
testowe.
Wyniki takich porównań rzadko przekładają się na odmienną konfigurację sprzętową,
a~inżynier stający przed wyborem algorytmu dla swojej platformy dysponuje
szczątkową i~trudno porównywalną wiedzą empiryczną.

\section{Cel i~zakres pracy}

Niniejsza praca ma na celu wypełnienie tej luki poprzez systematyczną ocenę
i~porównanie jedenastu algorytmów SLAM obejmujących wszystkie główne rodziny
metod i~modalności sensoryczne, przeprowadzone w~identycznych, w~pełni
kontrolowanych warunkach eksperymentalnych.

Badania przeprowadzono na robocie mobilnym TurtleBot3~Burger wyposażonym
w~trójwymiarowy skaner LiDAR, kamerę stereo, centralę inercyjną IMU oraz
enkodery odometrii kołowej --- platformie reprezentatywnej dla klasy małych
robotów mobilnych, a~jednocześnie dostatecznie bogatej sensorycznie,
by umożliwić testowanie wszystkich ocenianych rodzin algorytmów.
Oceniane metody obejmują fuzję sensoryczną, odometrię LiDAR, SLAM~2D i~3D
na bazie LiDAR-u, ciasno sprzężoną odometrię i~SLAM LiDAR-IMU oraz wizyjny SLAM.

Dane wejściowe stanowi jeden zbiór nagrań z~symulowanego środowiska magazynowego
w~NVIDIA~Isaac~Sim, który dostarcza dokładnych pozycji referencyjnych
eliminując potrzebę zewnętrznego systemu pomiarowego i~zapewniając pełną
powtarzalność warunków.
Wszystkie algorytmy uruchomiono pod wspólnym środowiskiem ROS\,2.

Ocena obejmuje dokładność trajektorii wyznaczaną miarami ATE i~RPE,
jakość map 3D mierzoną odległością Chamfera i~miarą F-score,
a~także koszt obliczeniowy wyrażony średnim obciążeniem procesora
i~szczytowym zużyciem pamięci RAM.

\section{Struktura pracy}

\begingroup
\setlength{\parskip}{2pt}
\setlength{\parindent}{1.5em}

Rozdział~2 zawiera podstawy teoretyczne niezbędne do zrozumienia omawianych
algorytmów, obejmując nawigację w~robotyce mobilnej, sensory, reprezentacje map,
rejestrację chmur punktów, filtracyjną i~grafową estymację stanu, preintegrację
IMU, podstawy wizyjnego SLAM oraz metryki ewaluacji.

Rozdział~3 opisuje każdy z~jedenastu ocenianych algorytmów pogrupowanych
według rodziny metod.

Rozdział~4 przedstawia konfigurację eksperymentalną, zbiór danych oraz wyniki
wszystkich jedenastu algorytmów w~zakresie dokładności trajektorii, jakości
mapy i~kosztu obliczeniowego.

Rozdział~5 podsumowuje wyniki eksperymentów, formułuje zalecenia dotyczące doboru
algorytmu SLAM do strukturyzowanych środowisk wewnętrznych oraz wskazuje kierunki
dalszych badań.

\endgroup
